% frontmatter.tex -- title, author, abstract, keywords, declarations (N01)
\title{Structure and computability of the single-quadratic existential fragment over the integers}
\author{Hiroki Fukui\\ \small Research Institute of Criminal Psychiatry / Sex Offender Medical Center, Tokyo, Japan\\ \small Department of Neuropsychiatry, Kyoto University, Kyoto, Japan\\ \small ORCID 0009-0008-7122-522X\\ \small ⟨EMAIL⟩}
\date{}
\maketitle
\linenumbers
\begin{abstract}
Let SE₂ be the class of relations on ℤᵏ definable by a \emph{single} polynomial equation of total degree at most two, with unrestricted integer unknowns.

Under a published input on ternary lattice cosets, f : ℤⁿ → ℤᵐ with two or more outputs is SE₂-definable exactly when f(a + Mt) = f(a) + Bt for a modulus M and an integer matrix B; a scalar function, exactly when it has this form, agrees with a rational polynomial of degree at most two, or is half the even one of two integer affine forms of opposite parity. The converses and the case without unknowns need no such input.

Finite fibres lie in a box of quadratic height and have polynomially bounded cardinality, with one constant depending on the equation alone; definable total functions therefore grow at most quadratically; they are computable; they are not closed under composition. SE₂-definable injections ℤⁿ → ℤᵐ satisfy n ≤ m, and this is sharp.

Unconditionally, with one constant C for all parameter points, each scalar fibre is bounded by C(1+‖a‖∞)² or carries a certificate pinning an infinite output subset of one of three forms: polynomial progressions, Pell-unit orbits, moving-modulus step orbits.

Two published results enter as named hypotheses: Grunewald–Segal decidability of quadratic solvability (code-uniform) and DPRM over ℕ. They give SE₂ ⊊ REC ⊊ RE = PP₂ = SE₄ at every positive arity. Sections 2–9 are formalized in Lean 4, as is Section 10 except Lemma 10.8, with the ternary input as an explicit hypothesis.

\end{abstract}
\noindent\textbf{Keywords.} Diophantine definability; existential fragments; quadratic equations; finite fibres; definable functions; computability.

\noindent\textbf{MSC 2020.} Primary 03D25; Secondary 11D09, 11U05, 03B35, 11Y50, 68V15.

\noindent\textbf{Funding.} This research did not receive any specific grant from funding agencies in the public, commercial, or not-for-profit sectors.

\noindent\textbf{Declaration of competing interest.} The author declares no known competing financial interests or personal relationships that could have appeared to influence the work reported in this paper.

